\chapter{Omega-3 Fatty Acids}

\medskip
\noindent\textit{\textbf{Under review.} This chapter is an AI-assisted educational draft; it carries no source citations and is not a substitute for professional medical advice.}

\section*{Definition and Overview}
Omega-3 fatty acids are polyunsaturated fatty acids characterised by a double bond at the third carbon from the methyl end of the chain. The three nutritionally and clinically important omega-3s are alpha-linolenic acid (ALA), eicosapentaenoic acid (EPA), and docosahexaenoic acid (DHA). ALA is an essential fatty acid found in plant sources such as flaxseed, walnuts, and chia, while EPA and DHA are predominantly found in oily fish and marine algae. EPA and DHA are incorporated into cell membranes and modulate inflammation, platelet function, and lipid metabolism. Dietary guidelines recommend regular consumption of oily fish, and fish-oil supplements (often containing icosapent ethyl, a purified EPA) are used for hypertriglyceridaemia. Interest also focuses on their role in cardiovascular disease, mental health, and inflammatory conditions.

\section*{Epidemiology}
In Australia, the National Health and Medical Research Council (NHMRC) recommends 1--2 serves of oily fish per week. Population surveys show most Australians do not meet these recommendations, with average seafood intake well below the guideline. Low omega-3 intake is more common among younger adults, those with lower socioeconomic status, and people who dislike fish. Globally, cardiovascular guidelines recommend omega-3 supplementation in people with elevated triglycerides who are at high cardiovascular risk. Use of fish-oil supplements is widespread, with a substantial proportion of Australian adults reporting regular omega-3 use. Rates of omega-3 deficiency are difficult to measure directly but are inferred from low dietary intake, and in Australia are not a common cause of clinically recognised deficiency disease because classic deficiency is rare.

\section*{Aetiology and Pathophysiology}
\begin{itemize}
    \item \textbf{Triglyceride lowering:} EPA and DHA reduce hepatic very-low-density-lipoprotein (VLDL) synthesis and secretion, lowering serum triglycerides in a dose-dependent manner
    \item \textbf{Anti-inflammatory effects:} compete with arachidonic acid for cyclo-oxygenase (COX) and lipoxygenase (LOX) enzymes, reducing pro-inflammatory eicosanoids and generating specialised pro-resolving mediators (resolvins, protectins)
    \item \textbf{Cell membrane structure:} DHA is highly concentrated in retinal and neuronal membranes, supporting membrane fluidity and function
    \item \textbf{Cardiovascular mechanisms:} effects on arrhythmia threshold, endothelial function, blood pressure, platelet aggregation, and possibly plaque stability
    \item \textbf{Mental health:} DHA/EPA incorporation into brain membranes may influence neurotransmitter signalling, inflammation, and neuroplasticity; mechanisms remain incompletely understood
    \item \textbf{ALA conversion:} only a small proportion of dietary ALA is converted endogenously to EPA ($\sim 5\%$) and DHA ($\sim 1\%$), so direct marine sources are more efficient at raising blood levels
\end{itemize}

\section*{Clinical Features}
Clinical manifestations relevant to omega-3 status:
\begin{itemize}
    \item \textbf{Deficiency is rare:} occurs mainly with severe malabsorption or very restricted diets; may cause dry scaly skin, poor growth in infants, and neurological symptoms in extreme cases
    \item \textbf{Hypertriglyceridaemia:} elevated fasting triglycerides; very high levels ($>$5.6 mmol/L) increase the risk of acute pancreatitis
    \item \textbf{Cardiovascular risk:} dyslipidaemia contributing to atherosclerosis and cardiovascular events
    \item \textbf{Supplement side-effects:} fishy aftertaste, reflux, nausea, bloating, and loose stools
    \item \textbf{Bleeding risk:} high-dose omega-3 may prolong bleeding time, particularly with anticoagulant or antiplatelet therapy
\end{itemize}

\section*{Differential Diagnosis}
Conditions and factors that overlap with omega-3-related presentations:
\begin{itemize}
    \item \textbf{Genetic dyslipidaemia:} familial combined hyperlipidaemia, familial hypertriglyceridaemia (type IV), familial chylomicronaemia (type I)
    \item \textbf{Secondary hypertriglyceridaemia:} diabetes, alcohol excess, hypothyroidism, renal disease, drugs (oestrogens, thiazides, beta-blockers, corticosteroids)
    \item \textbf{Atopic dermatitis (dry skin):} eczema, ichthyosis, nutritional deficiency of other fatty acids
    \item \textbf{Mood disorders:} major depression and anxiety (omega-3s studied as adjunctive therapy) -- require formal psychiatric diagnosis
    \item \textbf{Dyspepsia from supplements:} gastro-oesophageal reflux, peptic ulcer, functional dyspepsia
\end{itemize}

\section*{When to Seek Medical Care}
Most people do not need medical care specifically for omega-3 use. Seek medical review for persistent gastrointestinal symptoms on supplements, easy bruising or bleeding with high doses, or if considering omega-3 for cardiac risk reduction in combination with other medications.

\textbf{Call 000 for:}
\begin{itemize}
    \item Symptoms of acute pancreatitis (severe epigastric or upper abdominal pain radiating to the back, nausea, vomiting)
    \item Chest pain or other symptoms suggestive of a heart attack or stroke
    \item Unexpected or uncontrolled bleeding, especially on anticoagulants
    \item Severe allergic reaction to fish or fish-oil products (anaphylaxis) -- rare
\end{itemize}

\section*{Diagnosis and Investigations}
\begin{itemize}
    \item \textbf{Clinical/history:} dietary assessment of fish and omega-3 intake
    \item \textbf{Fasting lipid profile:} triglycerides, total cholesterol, LDL, HDL; repeat to confirm sustained hypertriglyceridaemia
    \item \textbf{Triglyceride thresholds:} normal $<$1.7 mmol/L; mild--moderate elevation 1.7--5.6 mmol/L; severe $>$5.6 mmol/L (pancreatitis risk)
    \item \textbf{Routine electrolyte/liver/thyroid:} to exclude secondary causes of hypertriglyceridaemia
    \item \textbf{Omega-3 index (research):} percentage of EPA+DHA in red-cell membranes as a marker of long-term intake (target typically 8--11\%)
    \item \textbf{Exclusion of pancreatitis:} serum lipase, amylase, abdominal imaging in acute presentation
\end{itemize}

\section*{Management}
\begin{itemize}
    \item \textbf{Dietary:} 2 serves of oily fish weekly (salmon, mackerel, sardines, tuna); plant sources for vegetarians (flaxseed, chia, walnuts, algae-based oils)
    \item \textbf{Supplementation for hypertriglyceridaemia:} icosapent ethyl (e.g., Vascepa) 2 g twice daily with food, or fish-oil EPA/DHA at doses of 2--4 g daily
    \item \textbf{Cardiovascular risk use:} icosapent ethyl reduced cardiovascular events in patients with elevated triglycerides on statin therapy (REDUCE-IT trial); discuss the evidence with a clinician
    \item \textbf{Dose and adherence:} take with meals to reduce reflux and fishy taste; up to 4 g daily is generally well tolerated
    \item \textbf{Combination therapy:} add statins or fibrates (e.g., fenofibrate) for persistent severe hypertriglyceridaemia
    \item \textbf{Monitor triglycerides} at 2--3 months after starting a therapeutic dose
    \item \textbf{Lifestyle:} reduce alcohol, refined carbohydrates, and simple sugars; regular exercise and weight management
    \item \textbf{Adjunct in psychiatry:} EPA-rich preparations (1--2 g daily) have been trialled alongside antidepressants in major depression, with modest evidence
\end{itemize}

\section*{Complications}
\begin{itemize}
    \item Acute pancreatitis with severe untreated hypertriglyceridaemia
    \item Increased bleeding risk when combined with anticoagulants or antiplatelet drugs
    \item Atrial fibrillation risk elevation reported in some high-dose trials (mainly high-dose EPA)
    \item Gastrointestinal intolerance (reflux, diarrhoea, bloating)
    \item Interaction with blood-pressure-lowering medication (additive hypotensive effect)
    \item Allergic reactions in individuals with fish or shellfish allergy
\end{itemize}

\section*{Prognosis and Outcomes}
In adults with hypertriglyceridaemia, omega-3 therapy typically lowers triglycerides by 20--30\% at doses of 2--4 g daily, with larger reductions at higher baseline levels. In the landmark REDUCE-IT trial, icosapent ethyl reduced the composite of cardiovascular events by approximately 25\% in high-risk statin-treated patients with elevated triglycerides. The broader evidence for omega-3 in primary prevention is mixed, and large meta-analyses show only modest or neutral cardiovascular benefit for low-dose fish oil in people without raised triglycerides. For most people, meeting dietary omega-3 recommendations is safe and associated with a well-balanced diet; supplementation at usual doses has a good safety profile. Outcomes in mental health remain variable, and omega-3s should be considered adjunctive rather than standalone therapy.

\section*{Prevention and Self-Care}
\begin{itemize}
    \item Eat oily fish 1--2 times weekly; choose low-mercury options in pregnancy (limit shark, swordfish, barramundi)
    \item Include plant omega-3 sources (flaxseed, walnuts, chia) for vegetarians and vegans
    \item Take fish-oil supplements with meals and refrigerate to minimise fishy aftertaste
    \item Check supplement labels for EPA/DHA content, not just ``fish oil''
    \item Have fasting lipids checked if there is a family history of high cholesterol or triglycerides
    \item Limit alcohol and simple carbohydrate intake to protect triglyceride levels
    \item Discuss supplements with a pharmacist or doctor if taking anticoagulants
\end{itemize}

\section*{Sources and Further Reading}
The following reputable sources provide further information for healthcare students and patients seeking reliable, current advice about omega-3 fatty acids.

\begin{itemize}
    \item National Health and Medical Research Council (NHMRC). \textit{Australian Dietary Guidelines.} https://www.nhmrc.gov.au/
    \item Heart Foundation (Australia). \textit{Omega-3 and heart health.} https://www.heartfoundation.org.au/
    \item Therapeutic Goods Administration (TGA). \textit{Regulation of fish oil complementary medicines.} https://www.tga.gov.au/
    \item National Heart, Lung, and Blood Institute (NHLBI). \textit{Cholesterol and triglycerides.} https://www.nhlbi.nih.gov/
    \item National Institutes of Health, Office of Dietary Supplements. \textit{Omega-3 fatty acids fact sheet.} https://ods.od.nih.gov/
    \item Healthdirect Australia. \textit{Omega-3 fatty acids.} https://www.healthdirect.gov.au/
\end{itemize}

\clearpage
